\chapter{Literature Review}
Reviews relevant existing research and demonstrates how your thesis fits into the broader academic conversation.

\section{Review of Previous Studies}

\subsection{Overview of Privacy-Preserving Mechanisms in Federated Learning}

\subsubsection*{Overall Landscape and Scope of Research}
The integration of Federated Learning (FL) into Intrusion Detection Systems (IDS) for IoT environments has resolved the issue of raw data centralization, yet the model weight sharing process itself still harbors data leakage risks~\cite{vyas2024}. To counter this, current studies classify Privacy-Preserving Federated Learning (PPFL) techniques into five main categories: Differential Privacy (DP), Homomorphic Encryption (HE), Secure Multi-Party Computation (SMPC), Trusted Execution Environments (TEE), and Blockchain~\cite{vyas2024}. Although TEEs provide a secure execution environment and Blockchain ensures transparency, they require specialized hardware or complex consensus mechanisms, which are unsuitable for the constrained resources of edge/IoT devices~\cite{manzano2025}. Therefore, this study implements DP and SMPC-based secure aggregation (SecAgg$^+$); HE is reviewed as an alternative but not implemented, because of its computational cost on constrained devices~\cite{nakayiza2026}.

\subsubsection*{Trusted Execution Environment (TEE)}
The Trusted Execution Environment (TEE) provides a completely isolated and secure hardware computing area for processing sensitive data, effectively preserving privacy. However, this technology is heavily dependent on specialized hardware security mechanisms and verifiable components. This strict hardware requirement makes TEE cumbersome and unsuitable for deployment across lightweight IoT networks comprising massively resource-constrained devices~\cite{vyas2024}.

\subsubsection*{Blockchain}
Blockchain preserves data privacy through decentralized ledgers and cryptographic algorithms that ensure tamper-resistance and eliminate the need for intermediaries. However, this system suffers from fatal performance weaknesses due to low throughput, high latency, and expensive verification costs stemming from continuously growing immutable storage. These severe latency and massive storage overheads make blockchain highly impractical for lightweight, multi-device IoT environments that demand rapid response times~\cite{vyas2024}.

\subsubsection*{Homomorphic Encryption (HE)}
HE enables computations to be performed directly on encrypted data, ensuring absolute security for model updates from clients before server aggregation without the need for decryption~\cite{nakayiza2026}. However, the computational cost and communication requirements of HE algorithms are extremely immense and heavy, leading to a surge in system latency~\cite{nakayiza2026}. This renders HE less feasible for systems requiring real-time responses and operating under resource constraints, such as IoT-based IDS~\cite{nakayiza2026}.

\subsubsection*{Differential Privacy (DP)}
Unlike the cumbersomeness of HE, Differential Privacy (DP) offers excellent scalability. The core principle of DP is to mask the contribution of any single data point by adding random noise to the model update at each client~\cite{ruzafa2022}. DP is exceptionally lightweight, consumes no extra bandwidth, and easily scales to tens of thousands of participating clients~\cite{ruzafa2022}. The tradeoff of this method is that injecting noise into the model inevitably degrades data accuracy, forcing engineers to constantly seek an optimal balance for the privacy budget~\cite{ruzafa2022}.

\subsubsection*{Secure Multi-Party Computation (SMPC/SecAgg)}
Conversely, Secure Aggregation (SecAgg) protocols based on SMPC comprehensively resolve the weakness of DP. SMPC prevents the server from directly observing individual client updates by ensuring that the server only receives the aggregated model update~\cite{bessaa2026}. Secure aggregation adds no statistical noise; with the fixed-point quantisation used by SecAgg$^+$, the aggregate differs from the plain sum only by a bounded quantisation error~\cite{bell2020}. As a trade-off, the original Secure Aggregation protocol requires multiple rounds of complex interactions between clients to exchange secret shares, causing severe bandwidth bottlenecks as the client population grows~\cite{bonawitz2017}; however, the subsequent SecAgg$^+$ protocol successfully resolves these overhead issues by replacing the complete communication graph with a $k$-regular graph, thereby reducing the per-client communication and computation complexity to a polylogarithmic scale~\cite{bell2020}.

\subsubsection*{Hybrid Privacy-Preserving Systems}
The study by Gutti et~al.~\cite{gutti2024} (integrating DP and SecAgg via the HADA algorithm) demonstrates that combining DP (to prevent inference attacks) and SecAgg (to protect the transmission channel) is a promising trajectory. However, stacking these security layers generates immense complexity in computation and hyperparameter tuning~\cite{gutti2024}. Therefore, before advancing to hybrid architectures, independently measuring the performance trade-offs between the DP pipeline and the SecAgg pipeline serves as a mandatory stepping stone to evaluate the endurance of the network system.

\subsubsection*{Distributed Differential Privacy}
Distributed DP combines the two families: each client adds only a share of the noise, and secure aggregation reveals only the sum, which carries the full calibrated noise~\cite{kairouz2021, agarwal2021}. Because secure aggregation operates on integers modulo a large number, discrete noise is preferred; the sum of independent Skellam variables is again Skellam, which permits exact privacy accounting after modular aggregation~\cite{agarwal2021}. The guarantee assumes that enough clients add their share honestly; with fewer honest clients the effective noise is smaller. Distributed DP fits releases that are a single sum, such as count tables, but not multi-step local DP-SGD, where each client's noise passes through non-linear local training before aggregation.

\subsection{Data Generation}

\subsubsection*{Overview of Data Generation Techniques}
In the field of cybersecurity and Intrusion Detection Systems (IDS) for Internet of Things (IoT) environments, Generative AI (GenAI) techniques for data generation are defined as deep generative paradigms capable of learning the underlying multivariate probability distribution of network traffic, thereby automatically synthesizing novel data samples whose structural attributes and statistical properties closely match genuine traffic distributions~\cite{deng2025}.

\noindent\textbf{Core problems addressed by data generation techniques:}
\begin{itemize}
    \item \textbf{Overcoming Extreme Class Imbalance:} Real-world IoT network traffic inherently exhibits extreme class skewness, where over 90--99\% of collected traffic consists of benign activity, while malicious behaviors---especially Low-Rate Stealth Attacks (LRSA) or rare DoS/DDoS variants---represent a negligible fraction. This acute shortage of attack samples severely biases downstream IDS classifiers toward majority classes, resulting in unacceptable false-negative rates on critical threats~\cite{deng2025}.
    \item \textbf{Beyond Interpolation-Based Oversampling:} Traditional statistical oversampling algorithms such as SMOTE~\cite{chawla2002} operate by performing linear interpolation between existing minority instances within a localized convex hull. As highlighted by Deng et~al.~\cite{deng2025}, deep generative models surpass these rigid boundaries by capturing the global multivariate distribution, enabling the synthesis of samples in unexplored regions of the feature space to cover edge cases and zero-day attack variants, thereby strengthening model robustness. Furthermore, because SMOTE requires simultaneous access to neighboring minority instances in a unified feature space, it is fundamentally ill-suited for federated environments where an individual edge gateway may hold few or no records of a rare attack class~\cite{chawla2002}.
    \item \textbf{Supporting Data Augmentation and Restructuring:} Systematic survey statistics indicate that up to 89\% of studies leverage GenAI specifically for data augmentation and class rebalancing, substantially improving the detection performance and generalization of IDS without exposing sensitive raw telemetry~\cite{deng2025}.
\end{itemize}

\noindent\textbf{Primary taxonomies of generative architectures in IDS:}
\begin{itemize}
    \item \textbf{Generative Adversarial Networks (GANs):} Architectures structured around an adversarial minimax game between a Generator network that creates synthetic samples from latent noise and a Discriminator network that evaluates sample authenticity. Due to their capacity to model sharp, high-dimensional distributions, GANs account for an overwhelming 58\% of research in generative IoT-IDS~\cite{deng2025}.
    \item \textbf{Autoencoders and Variational Autoencoders (AEs/VAEs):} Frameworks comprising an Encoder that maps high-dimensional traffic into a low-dimensional latent manifold and a Decoder that reconstructs the input. VAEs enforce continuous probability distribution priors over the latent space, permitting stable random sampling to synthesize realistic minority samples~\cite{deng2025}.
    \item \textbf{Transformer-Based Generative Models:} Architectures leveraging multi-head self-attention mechanisms to capture complex contextual correlations and sequential packet dependencies across temporal network traffic flows~\cite{deng2025}.
\end{itemize}

\subsubsection*{Analysis of Strengths and Weaknesses of Techniques in Existing Studies}
Although data generation techniques provide substantial potential for IoT-IDS, critical examination of existing literature reveals distinct computational bottlenecks, training instabilities, and privacy trade-offs when deployed in decentralized Federated Learning (FL) settings~\cite{deng2025}:
\begin{itemize}
    \item \textbf{Heavy Generative Denoising Techniques (Tabular Diffusion / Tab-DDPM):} Manzano et~al.~\cite{manzano2025} integrated tabular diffusion models (Tab-DDPM) directly at local FL clients to synthesize minority samples on the BoT-IoT benchmark. While Tab-DDPM achieves exceptional sample fidelity and distribution alignment, the reverse diffusion process requires iterative denoising over hundreds of sequential steps. On edge-tier testbeds (1 CPU core, 2\,GB RAM), generating a single synthetic sample requires approximately 32 seconds~\cite{manzano2025}. Such extreme computational latency is completely unviable for line-rate packet inspection in edge IoT environments.
    \item \textbf{Improved Policy-Guided GANs:} Yang et~al.~\cite{yang2025} employed an improved GAN guided by reinforcement learning to augment training sets for DDoS detection. While ablation experiments confirm that removing synthetic augmentation causes a noticeable drop in detection performance, the multi-component architecture requires up to 2.726 seconds per batch and consumes 19.6\,GB of GPU memory on the KDDcup99 dataset~\cite{yang2025}. Furthermore, the study did not isolate standalone generator inference latency, obscuring the true computational efficiency of the generator on edge hardware.
    \item \textbf{Survey Insights on Pure GAN Instability:} In their survey, Deng et~al.~\cite{deng2025} highlight that while GANs produce high-fidelity samples, the adversarial training dynamics are notoriously brittle and susceptible to mode collapse---wherein the generator repeatedly outputs a narrow set of duplicate patterns, wasting node compute without contributing meaningful diversity.
    \item \textbf{Random Resampling Techniques (Bootstrapping / Oversampling):} Jacob and Sultana~\cite{jacob2026} utilized bootstrapping (random resampling with replacement) at local nodes to balance the BoT-IoT and CSE-CIC-IDS2018 datasets. While computationally trivial and easily executable on microcontrollers, bootstrapping merely duplicates existing data points without expanding feature diversity. This encourages classifier overfitting and fails to generalize against evolving attack variants or concept drift~\cite{deng2025, jacob2026}.
    \item \textbf{Cloud-Based Centralized Data Generation Techniques (Cloud-Based VAE-GAN):} Anandhi and Srikanth~\cite{anandhi2026} proposed the SSAIM framework utilizing a hybrid VAE-GAN at the cloud coordinator to generate 51,200 Low-Rate Stealth Attack (LRSA) samples for the IoTID20 dataset, raising attack recall from 62\% to 98\%. However, placing data synthesis in the cloud violates the decentralized privacy principle of FL. Additionally, compressing the heavy VAE-GAN model via fixed-point quantization for transmission back to edge devices demonstrated that quantization below 6-bit (e.g., 4-bit) caused a severe degradation in detection sensitivity~\cite{anandhi2026}.
    \item \textbf{Federated Lightweight Edge GANs (TinyGAN in FL):} Durga et~al.~\cite{durga2026} deployed TinyGAN within a federated framework for IoMT malware detection, reporting 99.30\% accuracy after 20 rounds. However, to stabilize GAN training, all malware classes representing less than 0.5\% of the dataset were discarded, thereby sidestepping the core challenge of extreme class imbalance. Moreover, the study omitted standalone sample generation throughput measurements, leaving its actual computational advantage on constrained edge hardware unproven~\cite{durga2026}.
\end{itemize}

\section{Summary of the Literature Review}
\section{Contribution of Research}